% Clase 28 --- Suma de las dos ondas segun el desfasaje (§1.2--1.3).
% Tres casos con la MISMA geometria: lo unico que cambia es Delta(phi), que a su
% vez esta fijado por la diferencia de caminos. La amplitud resultante es
% 2A cos(Delta phi / 2).
\begin{center}
\begin{tikzpicture}

  \begin{axis}[
    fisfig, width=4.7cm, height=3.4cm, grid=none,
    xmin=0, xmax=6.5, ymin=-2.4, ymax=2.4,
    xtick=\empty, ytick=\empty, axis lines=middle,
    title={\small $\Delta\phi = 0$: refuerzo},
    title style={yshift=-2pt},
  ]
    \addplot[figgray!60, smooth, samples=120, domain=0:6.5] {sin(deg(x))};
    \addplot[figgray!60, smooth, samples=120, domain=0:6.5] {sin(deg(x))};
    \addplot[figblue, line width=1.2pt, smooth, samples=120, domain=0:6.5]
      {2*sin(deg(x))};
  \end{axis}

  \begin{scope}[xshift=5.0cm]
  \begin{axis}[
    fisfig, width=4.7cm, height=3.4cm, grid=none,
    xmin=0, xmax=6.5, ymin=-2.4, ymax=2.4,
    xtick=\empty, ytick=\empty, axis lines=middle,
    title={\small $\Delta\phi = \pi/2$}, title style={yshift=-2pt},
  ]
    \addplot[figgray!60, smooth, samples=120, domain=0:6.5] {sin(deg(x))};
    \addplot[figgray!60, smooth, samples=120, domain=0:6.5]
      {sin(deg(x)+90)};
    \addplot[figblue, line width=1.2pt, smooth, samples=120, domain=0:6.5]
      {sin(deg(x)) + sin(deg(x)+90)};
  \end{axis}
  \end{scope}

  \begin{scope}[xshift=10.0cm]
  \begin{axis}[
    fisfig, width=4.7cm, height=3.4cm, grid=none,
    xmin=0, xmax=6.5, ymin=-2.4, ymax=2.4,
    xtick=\empty, ytick=\empty, axis lines=middle,
    title={\small $\Delta\phi = \pi$: se cancelan},
    title style={yshift=-2pt},
  ]
    \addplot[figgray!60, smooth, samples=120, domain=0:6.5] {sin(deg(x))};
    \addplot[figgray!60, smooth, samples=120, domain=0:6.5]
      {sin(deg(x)+180)};
    \addplot[figblue, line width=1.2pt, samples=2, domain=0:6.5] {0};
  \end{axis}
  \end{scope}

\end{tikzpicture}\\[0.5em]
{\small\itshape Las dos ondas que llegan a $P$ tienen la misma amplitud y la
misma frecuencia; lo único que las distingue es la fase con que llegan, y esa
fase la fija la diferencia de caminos: $\Delta\phi = k(R_2-R_1)$. Sumar dos
senos de igual frecuencia da otro seno de la misma frecuencia,
\[
E = 2A\cos\!\Big(\tfrac{\Delta\phi}{2}\Big)\,
     \sen\!\Big(kR_1 - \omega t + \tfrac{\Delta\phi}{2}\Big),
\]
es decir que el desfasaje no cambia el ritmo de la oscilación sino
\textbf{su amplitud}, a través del factor $2A\cos(\Delta\phi/2)$. De ahí salen
las dos condiciones sin necesidad de más cuentas: amplitud máxima cuando
$\Delta\phi$ es múltiplo de $2\pi$ ---o sea, cuando la diferencia de caminos es
un número entero de longitudes de onda, $d\sen\theta = n\lambda$--- y
cancelación exacta cuando es un múltiplo impar de $\pi$, $d\sen\theta =
(n+\tfrac12)\lambda$. El caso intermedio del panel del medio no es ni una cosa
ni la otra: la amplitud queda entre $0$ y $2A$.}
\end{center}
